% -----------------------------------------------------------
\section{OT Extended Scripture Studies}
% -----------------------------------------------------------
	\label{EXSS-OT}
	
	% ----------------------
	\subsection{Commentary on the ``sons of Levi'' in Malachi 3:3}
	% ----------------------
		\label{EXSS-OT-MALACHI-3-3-SONS}
		
		Levi himself had three sons: Gershon, Kohath, and Merari (\hyperref[GENESIS-46-11]{Genesis 46:11}, \hyperref[EXODUS-6-16]{Exodus 6:16}). Gershon had two sons, Libni and Shimi (\hyperref[EXODUS-6-17]{Exodus 6:17}) and Merari had two, Mahali and Mushi (\hyperref[EXODUS-6-19]{Exodus 6:19}). Kohath had four sons: Amram, Izhar, Hebron, and Uzziel (\hyperref[EXODUS-6-18]{Exodus 6:18}). Then ``Amram took him Jochebed his father's sister to wife; and she bare him Aaron and Moses'' (\hyperref[EXODUS-6-20]{Exodus 6:20}). Going through the generations by the fathers (\hyperref[1CHRONICLES-6]{1 Chronicles 6}), then, you get: 
					
		\resizebox{.75\linewidth}{!}{%
		\Tree [.Levi [.Gershon ] [.Kohath [.Amram [.Aaron [.Eleazar ] [.Nadab ] [.Abihu ] [.Ithamar ] ] [.Moses ] [.Miriam ] ] Izhar Hebron Uzziel ] [.Merari ] ]
		}
					
		So Moses and Aaron are brothers, and they are the great-grandsons of Levi (though through their mother's side they skip a generation, because Amram's wife was Levi's daugther). All the descendants of Levi had special responsibilities over the tabernacle and all its implements (\hyperref[NUMBERS-3]{Numbers 3}, \hyperref[NUMBERS-4]{Numbers 4}, \hyperref[1CHRONICLES-23]{1 Chronicles 23--26}); they had special ministry and judgment responsibilities (\hyperref[DEUTERONOMY-21-5]{Deuteronomy 21:5}); they communicated the Lord's law to the people every seven years (\hyperref[DEUTERONOMY-31-9]{Deuteronomy 31:9}). They played a pretty fundamental role in the administration of ordinances and rituals among this early people.
		
		This verse appears to be the biblical basis for John says to JS and Cowdery in \hyperref[DC13]{DC 13}. See \hyperref[DEUTERONOMY-18-1]{Deuteronomy 18:1, 3}, \hyperref[DEUTERONOMY-26-4]{26:4}, and \hyperref[3NEPHI-24-3]{3 Nephi 24:3--4}.
		
		Here are some excerpts from commentaries:
		
		``Malachi had previously rebuked the priests in the second disputation for failing in their duties as ministers of the temple rituals and teachers of the Mosaic law (\hyperref[MALACHI-1-8]{1:8, 10, 13}; \hyperref[MALACHI-2-6]{2:6--8}). The priests are singled out as the target of God's refining process, `because they were the mediators between God and his people, and were therefore also responsible for the religious decline of the people. Thus purification of the people has to start with them.' Once purified, `they shall be fit to bring offerings to the LORD' (\textsc{NEB}; cf. \textsc{NLT}, `so that they may once again offer acceptable sacrifices to the LORD'). The prophet may have had the post-exodus Mosaic period in mind as the \textit{days gone by} or the ideal era of Hebrew covenant relationship with YHWH (cf. the reference to Phinehas in \hyperref[MALACHI-2-5]{2:5}). Taken together, the two expressions probably allude more generally to those times in Israelite history when the Hebrew priests and people faithfully served YHWH. In either case, the reference serves as an important reminder to the priests and the people of post-exilic Judah that they have a legacy they can aspire to emulate'' (\textit{Tyndale Old Testament Commentaries}, volume 28, \textit{Haggai, Zechariah, and Malachi}, Andrew E. Hill).
		
		``When the messenger of the covenant comes, he will proclaim a purifying and refining message, one especially directed to and effective among `the Levites,' that is, the priests (vv.2--4). Though the gospel narratives do not provide much direct evidence for this happening (but see, e.g., Nicodemus [\hyperref[JOHN-3-1]{John 3:1}, \hyperref[JOHN-19-39]{19:39}]; and many priests who later believed [\hyperref[ACTS-6-7]{Acts 6:7}]), there is no question that John's ministry paved the way for the life-changing message of the gospel and the establishment of the church (\hyperref[ACTS-19-1]{Acts 19:1--7})'' (\textit{The Expositor's Bible Commentary}, Revised Edition, Daniel--Malachi).
		
		``It is the Levites (lit., `the sons of Levi') who are purified. Levi has been mentioned twice earlier in Malachi in relation to a covenant of religious obligation with him (\hyperref[MALACHI-2-4]{2:4, 8}). His descendants, the Levitical priests (the terms `Levite' and `priests' are equivalent in Malachi), have reached a state of ritual impurity that disallows them from performing their expected duties. They therefore need to be made ritually pure...These Levitical priests, whose role was to purify the people through sacrifice, have so polluted the ritual practice and themselves that they now need purification before they can offer sacrifice on behalf of others. Those who raise themselves above the ordinances of God (\hyperref[MALACHI-1-7]{1:7--8}) now need the cleansing only achieved through their own abasement (\hyperref[MALACHI-2-9]{2:9}). When purification is complete, the Levites will be restored to their role as those who can effectively officiate in the sacrificial cult...The offering previously unacceptable from the hands of the priests (\hyperref[MALACHI-1-10]{1:10}; \hyperref[MALACHI-2-13]{2:13}) is now no longer so because of the righteousness in which it is brought...Yahweh did not previously condemn the Levitical priesthood as an institution, but rather the manner in which they had been performing their role. When the manner is rectified, the role can and must be resumed---not only for the sake of the priests but also for the people...The word `sacrifice' links verse 4 to the previous verse, now speaking specifically of the sacrifices of Judah and its capital, Jerusalem. Priestly sacrifices (v. 3) represent those of the people, but they must also precede them in time, since a defiled priest cannot perform acceptable rituals for a defiled people, but must purify himself first. The sacrifice of the people will then be `acceptable to the LORD' (\textsc{NIV}; `welcomed' \textsc{JB}; `pleasing/pleasant,' \textsc{KJV}, \textsc{NASB}, \textsc{NEB}, \textsc{NRSV}), a term used elsewhere of sacrifice (cf. those which are unacceptable, \hyperref[JEREMIAH-6-20]{Jeremiah 6:20}). This contrasts the previously unacceptable offerings (\hyperref[MALACHI-1-9]{Malachi 1:9--10}). Yahweh now accepts these offerings'' (\textit{The NIV Application Commentary: Joel, Obadiah, Malachi}, David W. Baker).
		
		``In the desert the Levites had been given the right to serve the ark. By virtue of this they became the keepers of the central sanctuary in the tribal league. It became their prerogative to expound the Mosaic law as well as sacrifice at the central sanctuary. As the Levites grew in number, not all of them could stay at the central sanctuary. This made some available for serving at local sanctuaries. Anyone could be a priest at a local sanctuary, but even there Levites were preferred. On the other hand, at the central sanctuary only the Levites could serve. These are the ``Levitical priests'' of Deuteronomy. All these Levites seem to be descendants of Moses and/or Ithamar. They were probably the Mushites and Libnites mentioned in \hyperref[NUMBERS-26-58]{Numbers 26:58a}. The Aaronites (Hebronites) and Korahites were not influential during this time, and we have no record of their activities'' (ABD 5281).
	
	% ----------------------
	\subsection{Comparisons of restoration-based versions of Malachi 4}
	% ----------------------
		\label{EXSS-OT-MALACHI-4-VERSIONS}
		
		Here are some tables comparing different versions of the verses in \hyperref[MALACHI-4]{Malachi 4} from different places.
		
		% -----
		\subsubsection{Malachi 4:1}
		% -----
		
			\begin{table}[H]
			\begin{tabular}{p{3cm}p{12cm}}
			\hyperref[MALACHI-4-1]{Malachi 4:1} & FOR, behold, the day cometh, that shall burn as an oven; and all the proud, yea, and all that do wickedly, shall be stubble: and the day that cometh shall burn them up, saith the LORD of hosts, that it shall leave them neither root nor branch.
			\end{tabular}
			\end{table}
			
		% -----
		\subsubsection{Malachi 4:2}
		% -----
		
			\begin{table}[H]
			\begin{tabular}{p{3cm}p{12cm}}
			\hyperref[MALACHI-4-2]{Malachi 4:2} & But unto you that fear my name shall the Sun of righteousness arise with healing in his wings; and ye shall go forth, and grow up as calves of the stall.
			\end{tabular}
			\end{table}
			
		% -----
		\subsubsection{Malachi 4:3}
		% -----
		
			\begin{table}[H]
			\begin{tabular}{p{3cm}p{12cm}}
			\hyperref[MALACHI-4-3]{Malachi 4:3} & And ye shall tread down the wicked; for they shall be ashes under the soles of your feet in the day that I shall do this, saith the LORD of hosts.
			\end{tabular}
			\end{table}
			
		% -----
		\subsubsection{Malachi 4:4}
		% -----
		
			\begin{table}[H]
			\begin{tabular}{p{3cm}p{12cm}}
			\hyperref[MALACHI-4-4]{Malachi 4:4} & Remember ye the law of Moses my servant, which I commanded unto him in Horeb for all Israel, with the statutes and judgments.
			\end{tabular}
			\end{table}
			
		% -----
		\subsubsection{Malachi 4:5}
		% -----
		
			\begin{table}[H]
			\begin{tabular}{p{3cm}p{12cm}}
			\hyperref[MALACHI-4-5]{Malachi 4:5} & Behold, I will send you Elijah the prophet before the coming of the great and dreadful day of the LORD: \\
			\hyperref[DC2-1]{DC 2:1} (1838/21 September 1823) & BEHOLD, I will reveal unto you the Priesthood, by the hand of Elijah the prophet, before the coming of the great and dreadful day of the Lord. \\
			\end{tabular}
			\end{table}
			
		% -----
		\subsubsection{Malachi 4:6}
		% -----
		
			\begin{table}[H]
			\begin{tabular}{p{3cm}p{12cm}}
			\hyperref[MALACHI-4-6]{Malachi 4:6} & And he shall turn the heart of the fathers to the children, and the heart of the children to their fathers, lest I come and smite the earth with a curse. \\
			\hyperref[DC27-9]{DC 27:9} (1835)\footnotemark & And also Elijah, unto whom I have committed the keys of the power of turning the hearts of the fathers to the children, and the hearts of the children to the fathers, that the whole earth may not be smitten with a curse; \\
			\hyperref[DC98-16]{DC 98:16--17} & Therefore, renounce war and proclaim peace, and seek diligently to turn the hearts of the children to their fathers, and the hearts of the fathers to the children; And again, the hearts of the Jews unto the prophets, and the prophets unto the Jews; lest I come and smite the whole earth with a curse, and all flesh be consumed before me. \\
			\hyperref[DC110-13]{DC 110:14--15} & Behold, the time has fully come, which was spoken of by the mouth of Malachi---testifying that he [Elijah] should be sent, before the great and dreadful day of the Lord come--- 15 To turn the hearts of the fathers to the children, and the children to the fathers, lest the whole earth be smitten with a curse--- \\
			\hyperref[DC2-2]{DC 2:2--3} (1838/21 September 1823) & And he shall plant in the hearts of the children the promises made to the fathers, and the hearts of the children shall turn to their fathers. 3 If it were not so, the whole earth would be utterly wasted at his coming. \\
			\hyperref[DC138-47]{DC 138:47--48} (1918) & The Prophet Elijah was to plant in the hearts of the children the promises made to their fathers, Foreshadowing the great work to be done in the temples of the Lord in the dispensation of the fulness of times, for the redemption of the dead, and the sealing of the children to their parents, lest the whole earth be smitten with a curse and utterly wasted at his coming.
			\end{tabular}
			\end{table}
				\footnotetext{%
					% \label{}%
					This part of DC 27 was not in the \hyperref[EMS-1830-JS-CIR-AUG-1830]{original revelation} given around August 1830; it was added later, to the 1835 DC. See JSR 72--80, reproduced \hyperref[EXSS-DC27-JSR]{here}, for discussion.
					}